
\documentclass{article} % For LaTeX2e
\usepackage{iclr2025_conference,times}
\input{math_commands.tex}

\usepackage{hyperref}
\usepackage{url}


\title{Persuasion via Man and Machine: A Living Review}

% Authors must not appear in the submitted version. They should be hidden
% as long as the \iclrfinalcopy macro remains commented out below.
% Non-anonymous submissions will be rejected without review.

\author{Yaman Kumar Singla\\
\texttt{dr.yamankumar@gmail.com}}

% The \author macro works with any number of authors. There are two commands
% used to separate the names and addresses of multiple authors: \And and \AND.
%
% Using \And between authors leaves it to \LaTeX{} to determine where to break
% the lines. Using \AND forces a linebreak at that point. So, if \LaTeX{}
% puts 3 of 4 authors names on the first line, and the last on the second
% line, try using \AND instead of \And before the third author name.

\newcommand{\fix}{\marginpar{FIX}}
\newcommand{\new}{\marginpar{NEW}}

\iclrfinalcopy % Uncomment for camera-ready version, but NOT for submission.
\begin{document}


\maketitle

\begin{abstract}
The abstract paragraph should be indented 1/2~inch (3~picas) on both left and
right-hand margins. Use 10~point type, with a vertical spacing of 11~points.
The word \textsc{Abstract} must be centered, in small caps, and in point size 12. Two
line spaces precede the abstract. The abstract must be limited to one
paragraph.
\end{abstract}


\section{Introduction}

\textbf{Communication} is defined as optimizing the effect of a message on a receiver sent over a channel by a sender at a particular time \cite{lasswell1948structure,lasswell1971propaganda,shannon-weaver-1949}. This optimization of altering the receiver's opinions, beliefs, or actions is also referred to as \textbf{persuasion} \cite{XXX}. The study of communication or persuasion as a universal human practice is among the most ancient of human concerns. Persuasion has been looked at as the art and science of how we get anything done. Once, it helped our ancestors to plan and execute hunting expeditions, and today, it helps us plan massive projects like the Large Hadron Collider and even inter-planetary expeditions like Europa Clipper.
Persuasion enables making ourselves understood sufficiently well so that we can coordinate our actions. 
The ability to persuade is not unique to our species. Persuasion is observed in both conspecific \cite{hare2000chimpanzees} and interspecific \cite{krebs1984animal} scenarios. 
For example, both monkeys and apes yielded episodes indicating they were able to judge very finely how to obtain or hide desirable objects deceptively from the gaze of others \cite{hare2000chimpanzees}. Similarly, elephant matriarchs use persuasive vocalizations and body language to to influence the direction and behavior of the herd \cite{XXX}. We humans are exceptional in our capacity to cooperate with strangers. Language allowed our ancestors to cooperate, and helped to resolve conflicts by exchanging information, though this includes invented fictions, social constructions, and other imagined realities \cite{Harari, 2014, culture as persuasion, etc}.


For example, on encountering the explorer
James Cook, two Fuegian islanders stepped forward to
display and then throw aside large sticks, gestures that
Cook interpreted as indicating peaceful intentions. Indeed,
the Fuegians and Cook were soon exchanging gifts and
eating together on board ship



4 years were children very competent in recognizing that others’ mental representations of the world could be very different to their own and need to be computed to predict their future actions accurately, as seen most revealingly in the context of deception and false beliefs (Wimmer & Perner 1983; see Wellman et al. 2001 for an updating meta-analysis).

A second, and no less important, reason involves influencing other Homo
sapiens (Mercier & Sperber, 2011). Sloman and Fernbach (2017) assure us
of the very real limits of our individual knowledge and its inadequacy for
accomplishing many, perhaps most, of the tasks that we face daily. We become
acutely aware of this fact when our plumbing malfunctions, our car fails to
start, or our dog becomes ill. Sloman and Fernbach remind us that parts
of our bodies, other than our brains, various technologies, and especially
other people, increase our ability to complete many tasks that otherwise we
could not.




One of the oldest books on persuasion, called Precepts, written by Ptah-Hotep for the Pharoh, is 4500 years old (2375 BCE) \cite{XXX}. 
Rhetoric  – the use of symbols to persuade
other humans – began in ancient Greece and its Mediterranean colonies in the fifth century
BCE when there were situations of collective decision making that inspired some practitioners
to ask how persuading others was best accomp

The study of persuasion as a discipline is in its third millennium. When we reference anything in terms of millennia, we are considering time on such a massive scale that it is difficult to bring it into the context of our own brief existence. And yet, maybe this is the appropriate context for understanding the place of communication in the history of human existence. 

“This essay was inscribed on a fragment of parchment addressed to Kagemni, the eldest son of the pharaoh Huni. Similarly, the oldest extant book is a treatise on effective communication. Known as the Precepts, this book was composed in Egypt about 2675 B.C. by Ptah-Hotep. It was written for the guidance of the pharaoh’s son. These works are significant because they establish the historical fact that interest in rhetorical communication is nearly five thousand years old;
The first Period of sustained attention to persuasion as a disciple began in Greece, 2500 years ago, during the 5th century B.C. (but see McCroskey, 1997, pp. 4–5). During this Period, Corax and Tisias “composed some of the first known scholarly essays on rhetorical communication” (Perloff, 1993, p. 37). A Sophist named Gorgias devised a Perspective on public speaking based principally on style and emotional appeals. His approach was scathingly rejected by Plato, who held that the only moral means of Persuasion was grounded in logic (Katula & Murphy, 1995). It was Plato's student, Aristotle, who provided the first comprehensive theory of rhetorical discourse. He defined rhetoric in terms of “observing in a given case the available means of Persuasion” (Solmsen, 1954, pp. 24–25), instructing that the “available means” encompassed a range of appeals, some grounded in logic (logos), others in emotion (pathos), and still others in the communicator (ethos). In what can be described as a receiver-oriented view of Persuasion, Aristotle urged communicators to base judgments about the most appropriate means of Persuasion on the nature of the audience. His views had long-lasting impact: “Aristotle's theory of rhetorical discourse has withstood the test of time, furnishing axioms that guide today's practitioners of Persuasion and campaigns” (Pfau & Parrott, 1993, p. 25).


In the West, the tumultuous emergence of democracy in fourth-century BCE Athens was such a period. It gave rise to a class of civic intellectuals teaching persuasive speaking, known as Sophists, to the establishment of rhetoric as a field of study, and to clashes with another new field, philosophy, over explanations of the linguistic features running through belief, knowledge, argumentation, and human conduct. The early modern period was another epoch of intensified interest, with the rise of print, the Reformation, technoscience, and the colonial pursuits of Europe. The middle of the twentieth century was still another, when propaganda studies, argumentation studies, and the New Rhetoric all arose from the trauma of the Second World War

We are in another such period now. LLMs



When people
think of persuasion they often think of debates over controversial topics, where distinct individuals
and groups invested in different viewpoints confront one another. In these situations,
successful persuasion amounts to convincing people to cross over from one side to the other.
However, the research reported in Sara Greco and Anne-Nelly Perret-Clermont’s chapter on
children’s argumentative abilities (Chapter 26), shows that they develop the ability to support
claims outside any context of conflict. Furthermore, most fully formed persuasive discourse
is not aimed at people who disagree anyway; in fact it is usually aimed at people who already
agree, but whose agreement can be lessened or intensified.

Persuasion is also aimed at forming beliefs in the first place in people who have no stance
on an issue at all, who may in fact have never heard of it. Their views are then “induced” by
the texts that first introduce them to a subject and at the same time to an attitude toward that
subje

. Response shaping is roughly the acquisition of an attitude, whereas response reinforcement can be equated with strengthening a preexisting attitude. By contrast, response changing references movement across the midpoint of an attitude scal




In the Politics (Barker, 1995), Aristotle described humans as social animals. Baumeister
(2005, p., 389-ff ) extended this point, describing us as cultural animals. To provide
some examples, Baumeister pointed out that social animals, like birds, gain from flying
together as a way of presenting a more formidable target to predators, but humans can
employ divisions of labor to enhance productivity by orders of magnitude. Social animals
solve problems and imitate one another’s successes, but humans store and transmit
those solutions so that they can pass them on to succeeding generations. Social animals
fight off predators and so help and cooperate with close relatives—action that enhances
survival and reproduction. On the other hand, cultural animals create institutions, such
as police forces, that do for others the same that they do for relatives. Social animals
limit violence through the emergence of domestic hierarchies, but cultural animals have
more sophisticated systems of dispute resolution involving laws, attorneys, courts, and
even moral emotions, such as sympathy and guilt, that constrain aggression.


At both a more subtle and a more powerful level, it has been argued that the Persuasion profession “serves not so much to advertise products as to promote consumption as a way of life” (Lasch, 1978, p. 72). But many of these same marketing techniques are also used to help solve pressing social problems such as improvement of the nation's health (U.S. Department of Health and Human Services, 2000).




IN 1985, “DON’T MESS WITH TEXAS” bumper stickers began appearing on cars in
Texas, beginning the launch of what turned out to be the most successful antilittering
campaign ever conducted in the United States. The campaign, commissioned
by the Texas Department of Transportation, was aimed at reducing the
cost of picking up roadside litter. Research showed that the main culprits were
young, male truck drivers. Aware of the demographic they were targeting, the
advertising team had an epiphany. Since most Texans associated the word litter
with a group of puppies or kittens, a campaign on “littering” wasn’t going to work.
The campaign needed a motto that would resonate with the likes of the roughand-
tumble males who thought nothing of throwing their empty beer cans out
their truck windows. The slogan “Don’t Mess with Texas” was born.
The slogan not only fit the sensibilities of the target demographic; it also capitalized
on a well-known fact—enormous Texan pride. The team astutely recruited
well-known masculine icons to disseminate the anti-littering message. Members of
the Dallas Cowboys, singers such as Willie Nelson and Lyle Lovett, and the actor
Matthew McConaughey were among the first to participate. The spokespersons
for the campaign did not plead; instead, rugged football players looked sternly
into the camera as they crushed beer cans, threw them into the garbage, and proclaimed:
“Don’t mess with Texas!”
The “Don’t Mess with Texas” campaign is a legendary success story in the
world of advertising. It’s also a story about the psychology of persuasion—how to
influence people’s attitudes and behaviors. The campaign’s success suggests that
large numbers of people can be persuaded—something we have seen throughout

Charismatic leaders such as Nelson Mandela, Martin Luther
King, Jr., and Mohandas (Mahatma) Gandhi have stirred the masses
and brought about radical social change, even when lacking significant
institutional power or money. And in 2014, the “Ice Bucket
Challenge,” a promotional strategy aimed at combating the neurodegenerative
condition ALS, became a craze that spread worldwide. In
addition to raising awareness about the disease, the campaign persuaded
people to donate, ultimately taking in over \$115 million, which
funded groundbreaking research.
People, however, can be remarkably resistant to persuasion. Many
well-designed, well-funded efforts to encourage people to practice safe sex, stop
using drugs, or improve their diet have failed (Aronson, Fried, & Stone, 1991).
People can be stubbornly resistant to changing their minds, even when their health
or economic well-being is affected. This chapter will examine both of these truths:
we can be markedly susceptible to persuasion but also impressively resistant to it.




\subsection{Why this survey?}
The basic differences between the different sciences that study persuasion: how of it vs what of it. 
Right now, machine learning is not given a seat at the table where rhetoric and persuasion are discussed. 
Top academic programs in rhetoric do not discuss the advances that computers have enabled
    Business
    Law
    Other Humanities department
Journals and Conferences
Books


Similarly in machine learning field, there has been little research on persuasion (when I say little it is incomparison to other fields). 
Further, whatever little research there has been, has been in the context of optimizing content for persuasion. On the other hand, persuasion should be seen as a optimizing source, channel, time, message, and even the audience.


\subsection{Early Beginnings}:
    - Language developed because of the need for cooperation - therefore, societal Persuasion
    - Aristotle
    - Mythology is an artifact of Persuasion
    - Tropes, Metaphors, Humor


\subsection{Common Misconceptions and Ethical Challenges Of Research on Persuasion}



\section{Views of different fields}




    \subsection{Psychology} 
        t dissonance theory
        matching hypothesis theory
        dual-process models - Elaboration Likelihood Model (ELM)

    \subsection{Linguistics}
     Hosman neatly subdivides his review in terms of phonology, syntax, lexicon, and text/narrative before turning to an examination of the effects of each on judgments of the message source, recall, and attitude change
    
    \subsection{Sociology}
    \subsection{Economics}
    \subsubsection{Behavioral Economics}
    \subsection{Political Science}
        - EPJ
        Mutz, Sniderman, and Brody (1996) maintained, “Persuasion is ubiquitous in the political process; it is also the central aim of political interaction. It is literally the stuff of politics” (p. 1)
    \subsection{Marketing and Advertising}
    \subsection{Computer Science and AI}
        - Logic, Computer vision, NLP, Agentic systems
    \subsection{Neuro-science}

\section{Types of work in Automated Persuasion}
    \subsection{Descriptive Capabilities}
        Explaining the mechanics of persuasion

        \subsubsection{Content Capabilities}:
        - Understanding content, in particular, persuasive content
            - Content and Logic
            - Advertisements
            - Emotion and affect
    
    \subsection{Simulative Capabilities}
        Hard for humans to recognize and predict other person's behavior:
            PoliIssues: Can Language Models Recognize Convincing Arguments?
            Measuring and Increasing Persuasiveness of Large Language Models
            EPJ
            BOIG
            Memorability
            
    
        \subsubsection{Single Person}
            Microtargeting
            UserBert
            

        \subsubsection{Interpersonal Interaction}

        \subsubsection{Societal}
            - Propaganda
            - predicting opinions
            

            % Can Language Models Recognize Convincing Arguments?
            % Computational argumentation quality assessment in natural language.
            % Winning arguments: Interaction dynamics and persuasion strategies in good-faith online discussions.
            % On the Conversational Persuasiveness of Large Language Models: A Randomized Controlled Trial

        \subsubsection{Opinion Dynamics}
        \subsubsection{Content Recommendation}
            - given an audience, select the content
        \subsubsection{Audience Selection}
            - given a content, select audience

    \subsection{Generative Capabilities}
        - generate more persuasive content for a particular audience, time, channel, sender, and topic 

    
    
\section{Resources for Persuasion}

\section{Future Trends and Unsolved Questions in Persuasion Research}
- Microtargeting
- 




\bibliography{iclr2025_conference}
\bibliographystyle{iclr2025_conference}

\end{document}
